\section*{Chapter \ref{ch:m2:the-rates-options-market} --- The Rates Options Market}

\begin{solution}{exo:m2:the-rates-options-market:1}
$95/\sqrt{252} = 5.98$ basis points a day; $7\sqrt{252} = 111.1$ basis points a
year.
\end{solution}

\begin{solution}{exo:m2:the-rates-options-market:2}
The company fears higher rates: it buys a \omterm{def:m2:the-rates-options-market:cap}{cap} on its loan's rate, or a \omterm{def:m2:the-rates-options-market:swaption}{payer
swaption} if it wants the right to fix its rate later. The pension fund fears
lower rates: it buys a \omterm{def:m2:the-rates-options-market:swaption}{receiver swaption}, or a \omterm{def:m2:the-rates-options-market:cap}{floor} on floating income.
\end{solution}

\begin{solution}{exo:m2:the-rates-options-market:3}
Black models the rate as lognormal, always positive; a negative forward or
strike has no logarithm. \omterm{def:m2:the-rates-options-market:normal}{Normal volatility} handles any sign without a free
parameter, is quoted in the basis points traders already think in, and its
daily equivalent compares directly with realised moves; a shifted lognormal
model works too, but its quotes depend on the shift, which must be agreed.
\end{solution}

\begin{solution}{exo:m2:the-rates-options-market:4}
Half the straddle: $7.80 \times 0.0095/\sqrt{2\pi}$ of notional, USD~2.96
million. The straddle's break-even is 75.8 basis points either way.
\end{solution}

\begin{solution}{exo:m2:the-rates-options-market:5}
15.0\%, 45.4\% and 326.7\%.
\end{solution}

\begin{solution}{exo:m2:the-rates-options-market:6}
In the normal model the rate at fixing is symmetric around the forward, 4\%,
so the chance and size of ending 50 basis points above are the same as 50 below:
\omterm{def:m2:the-rates-options-market:cap}{caplet} and floorlet have equal values period by period. In Black's model the
distribution of the rate is skewed to the right, and at the same Black
volatility the \omterm{def:m2:the-rates-options-market:cap}{floor} at 3.5\% would be worth less than the \omterm{def:m2:the-rates-options-market:cap}{cap} at 4.5\%.
\end{solution}

\begin{solution}{exo:m2:the-rates-options-market:7}
0.359\%: the at-the-money normal price per unit annuity is
$0.009/\sqrt{2\pi} = 0.00359$, and a Black call is worth less than the forward,
so no Black volatility reproduces it below that forward.
\end{solution}

\begin{solution}{exo:m2:the-rates-options-market:8}
Lognormal volatility is a percentage of the rate's level. At 20\% on a 5\% rate
the market expects moves of about 100 basis points a year; at 60\% on a 1\% rate,
about 60. In basis points, the units in which a position gains or loses, the
second currency is the calmer one. Compare normal volatilities.
\end{solution}

\begin{solution}{pb:m2:the-rates-options-market:1}
\textbf{1.} A \omterm{def:m2:the-rates-options-market:swaption}{receiver swaption}, 3y$\times$7y struck at 4.5\%: calling the bond in
three years stops its 4.5\% payments for seven years, which is worth having
exactly when the seven-year rate is below 4.5\%.
\textbf{2.} Cancelling a swap in which it pays 4.5\% is the same as then entering
one in which it receives 4.5\%: a \omterm{def:m2:the-rates-options-market:swaption}{receiver swaption} on the same terms.
\textbf{3.} 4.00\%; annuity 5.336.
\textbf{4.} In the money by 50 basis points: a receiver struck at 4.5\% on a 4\%
forward.
\textbf{5.} USD~25.87 million, 517 basis points of notional.
\textbf{6.} USD~176\,826 per basis point of \omterm{def:m2:the-rates-options-market:normal}{normal volatility}.
\textbf{7.} $-\text{USD}~163\,698$ per basis point: the receiver gains when rates
fall.
\textbf{8.} USD~353\,651.
\textbf{9.} USD~73\,740.
\textbf{10.} About USD~240 million ($176\,826/73\,740 \times 100$ million).
\textbf{11.} In the swap's terms: the dealer charges it a lower floating spread,
or pays it an upfront amount, for the cancellation right.
\textbf{12.} 63.8 basis points a year: USD~25.87 million over the ten-year
annuity of 8.11 on USD~500 million.
\textbf{13.} When refinancing for seven years costs it less than 4.5\%, roughly
when the seven-year rate plus its credit spread is below the coupon; the dealer
cancels the swap on the same day, and the bank is left with neither.
\textbf{14.} Investors in the callable are paid for the option in a higher
coupon; the bank sells the same option to the dealer and pays floating like any
other issuer. If the dealer pays more for the option than the investors
charged, the bank's floating cost is lower than on a plain bond swapped plainly;
it also reaches investors who want the extra coupon.
\textbf{15.} The receiver moves deeper into the money: its value rises to
USD~34.76 million, its delta grows, its vega falls to USD~156\,048; the dealer
must rebalance both its delta and its vega hedges.
\textbf{16.} Dealers sell the volatility they acquire, pushing long-dated
volatility down, most in the expiries and tenors of the calls.
\textbf{17.} Investors who need options: mortgage investors hedging negative
convexity, insurers and pension funds hedging guarantees, and speculators who
think volatility cheap.
\textbf{18.} Callable every year, it is a Bermudan \omterm{def:m2:the-rates-options-market:swaption}{swaption}: the bank will call on
the best date, which depends on the whole future curve; the value and the
hedge need a model of the cube's dynamics, and the hedge is a strip of
European \omterm{def:m2:the-rates-options-market:swaption}{swaptions} whose weights change with the market.
\textbf{19.} Named result: \emph{the callable's vega} is USD~176\,826 per basis
point of \omterm{def:m2:the-rates-options-market:normal}{normal volatility}, which the dealer lays off by selling about USD~240
million of 3y$\times$7y at-the-money straddles.
\textbf{20.} From issuers of callable debt, who sell the option in their bonds to
dealers, who sell it on.
\end{solution}

\begin{solution}{iq:m2:the-rates-options-market:1}
Lognormal volatility measures moves in proportion to the rate's level; \omterm{def:m2:the-rates-options-market:normal}{normal
volatility} in basis points, independent of the level. They agree at the money
up to a factor of the forward. It matters when rates are low or negative,
where lognormal quotes explode or do not exist, for smiles, which the two
models shape differently, and for hedging: the delta of an option differs
between the models.
\iqlookfor{units, the low-rate breakdown, and the model dependence of the smile
and the delta.}
\end{solution}

\begin{solution}{iq:m2:the-rates-options-market:2}
A payer is the right to enter a swap paying fixed at the strike, valuable if
rates rise; a receiver the right to receive fixed. Borrowers planning to fix
their debt, and investors hedging rate rises, buy payers; investors hedging
falling rates, mortgage investors hedging prepayments and pension funds buy
receivers; callable issuers, through their dealers, sell receivers.
\iqlookfor{payoffs, and the natural buyers and sellers.}
\end{solution}

\begin{solution}{iq:m2:the-rates-options-market:3}
With $F_T = F + \sigma\sqrt T Z$ and $K = F$, the payer is worth
$A\,\mathbb{E}[\sigma\sqrt T Z^+] = A\sigma\sqrt T\,\mathbb{E}[Z^+] =
A\sigma\sqrt T/\sqrt{2\pi}$, since $\mathbb{E}[Z^+] = \varphi(0) = 1/\sqrt{2\pi}$.
The straddle is twice that, about $0.8\,A\sigma\sqrt T$.
\iqlookfor{the one-line derivation and the $0.4$ rule of thumb.}
\end{solution}

\begin{solution}{iq:m2:the-rates-options-market:4}
Each day the hedged straddle earns roughly its gamma times the squared move
over two, and pays its time decay; the two balance when the day's move equals
the implied daily volatility, 95 over $\sqrt{252}$, about 6 basis points. Days
with larger moves make money, smaller lose; compare the realised daily moves,
weighted by gamma, with six.
\iqlookfor{gamma against theta, and the daily break-even move.}
\end{solution}

\begin{solution}{iq:m2:the-rates-options-market:5}
Callable issuers sell their call options to dealers inside swaps; dealers sell
the volatility on, concentrated in long expiries and tenors, which pushes
their volatility down. Mortgage investors are short prepayment options and hedge
by buying \omterm{def:m2:the-rates-options-market:swaption}{swaptions}, or by trading dynamically in the direction of moves, which
raises realised and implied volatility.
\iqlookfor{the two structural flows and their opposite signs.}
\end{solution}

\begin{solution}{iq:m2:the-rates-options-market:6}
Use a closed-form or rational approximation of the inverse Bachelier function
(accurate to near machine precision, as published for this purpose) rather
than a root finder, vectorised over options; handle the edge cases (prices at
or below intrinsic value, zero time, deep wings) explicitly. Test by round trips
on a grid of moneyness and maturities, against a slow bisection reference, on
random inputs, and on the edge cases; benchmark throughput on production-like
batches.
\iqlookfor{an explicit inverse, edge cases, and round-trip testing.}
\end{solution}
